% !TEX root = hw5-answers.tex
\chead{\textbf{Exercises week 5}}

\begin{document}
\flushleft\textbf{Topic: Markov kernels, Conditional Independence}
\section{Conditional Independence of Random Variables}
\begin{enumerate}
\item Consider the discrete random vector $(X, Y, Z)$ with joint PMF given in table 1.

\begin{table}[htb]
\centering
\begin{tabular}{llll}
\cline{3-4}
                                           & \multicolumn{1}{l|}{}    & \multicolumn{1}{l|}{y=1} & \multicolumn{1}{l|}{y=2} \\ \hline
\multicolumn{1}{|c|}{\multirow{3}{*}{z=1}} & \multicolumn{1}{l|}{x=1} & \multicolumn{1}{l|}{2}  & \multicolumn{1}{l|}{4}  \\ \cline{2-4} 
\multicolumn{1}{|c|}{}                     & \multicolumn{1}{l|}{x=2} & \multicolumn{1}{l|}{3}  & \multicolumn{1}{l|}{6}  \\ \cline{2-4} 
\multicolumn{1}{|c|}{}                     & \multicolumn{1}{l|}{x=3} & \multicolumn{1}{l|}{1}  & \multicolumn{1}{l|}{2}  \\ \hline
\multicolumn{1}{|l|}{\multirow{3}{*}{z=2}} & \multicolumn{1}{l|}{x=1} & \multicolumn{1}{l|}{12}  & \multicolumn{1}{l|}{12}  \\ \cline{2-4} 
\multicolumn{1}{|l|}{}                     & \multicolumn{1}{l|}{x=2} & \multicolumn{1}{l|}{6}  & \multicolumn{1}{l|}{6}    \\ \cline{2-4} 
\multicolumn{1}{|l|}{}                     & \multicolumn{1}{l|}{x=3} & \multicolumn{1}{l|}{0}    & \multicolumn{1}{l|}{0}  \\ \hline
\end{tabular}
\caption{Unnormalized joint probability mass function p(x, y, z).}
\end{table}\label{discrete-ci}

Check if $X$ is conditionally independent of $Y$ given $Z$. Provide the steps through your derivations. (1pt)
\begin{gradedsolution}
% The matrix rank of both $3\times 2$ of unnormalized $P(XY|Z)$ is 1, implying that it arises from an outer product $p(x,y|z)=p(x|z)p(y|z)$.
% If non-negative matrix $M$ can be factorised as $M=vw^T$, then either $v$ and $w$ are non-negative or $v$ and $w$ are non-positive (as the outer product would otherwise have a negative number). Consequently, taking if necessary $(v,w) \mapsto (-v, -w)$, we can choose a factorization into non-negative $v,w$.
Given a rank one matrix $M$ with factorization $M=vw^T$, all other factorizations are related to this factorization by $M=(v \lambda)(w / \lambda)^T$ for $\lambda \in \mathbb{R}\setminus\{0\}$. Choose $\lambda$ such that $\sum v_i=1$. Then if $M$ is normalized, $1=\sum_{ij}M_{ij}=\sum_{ij}v_iw_j=\sum_jw_j$, $w$ is also normalized. Consequently, $v_i=v_i \sum_j w_j = \sum_j M_{ij}$ and $w_j= w_j \sum_i v_i = \sum_i M_{ij} $. If $M$ is non-negative, then $v$ and $w$ are also non-negative. Hence, $v$ and $w$ are the marginals of $M$.

Hence, for any joint distribution $P(X, Y|Z=z)$, we only need to check if the matrix $p(x, y|z)$ is rank-one, as it can then be written as $p(x,y|z)=p(x|z)p(y|z)$. This is here the case.

Alternatively, it's easy to read off $p(x|z)$ and $p(y|z)$ from the joint distribution and verify $p(x,y|z)=p(x|z)p(y|z)$.
\end{gradedsolution}

\item Let $X,Y,Z,U$ be discrete random variables (on discrete spaces) with joint probability mass function $p(x,y,z,u)$. Show that the following 6 separoid axioms for conditional independence hold:

\begin{enumerate}
\item Redundancy: $X \Indep Y \;|\; X$

\item Symmetry: $X \Indep Y\; |\; Z \implies Y \Indep X\; |\; Z$ 

\item (Left) Decomposition: $X, U \Indep Y \; | \; Z \implies U \Indep Y | Z\; $

\item (Left) Weak Union: $X,U \;\Indep \;  Y \; | \; Z \implies  \; X \Indep Y \; |\; U,Z $

\item (Left) Contraction: $(X \Indep Y\; |\; U, Z) \; \land \; (U \Indep Y | Z) \implies X, U \Indep Y \; | \; Z$ \\

\item Intersection: For the following part, assume that p is strictly positive, i.e. for all values $x,y,z,u$ we have that $p(x, y, z, u) > 0$. Then show that also the Intersection property holds:
$$(X \Indep Y \;|\;U,Z) \land (X \Indep U \; | \; Y,Z)  \implies  X \Indep U,Y \; | \; Z.$$
\end{enumerate}
(6$\times$1 pts)\\
\end{enumerate}

\begin{gradedsolution}
  $X \Indep Y \; | \; Z$ is true w.r.t. a PMF $p(x, y, z)$ if a PMF $q(x|z)$ exists, such that for all $x \in \S X, y \in \S Y, z \in \S Z$:
\[
  p(x,y,z)=q(x|z)p(y,z)
\]
where $p(y,z)$ is the marginal of $p(x,y,z)$.
\begin{enumerate}
\item Redundancy: $X \Indep Y \;|\; X$

\begin{align*}
p(x_1,y,x_2)&=\delta_{x_1x_2}p(y, x_2)\\
q(x_1|x_2)&:=\delta_{x_1x_2}
\end{align*}

\item Symmetry: $X \Indep Y\; |\; Z \implies Y \Indep X\; |\; Z$ 

We are given a $q(x|z)$ s.t. $p(x,y,z)=q(x|z)p(y,z)$. By marginalisation over $Y$: $p(x, z)=q(x|z)p(z)$. By conditioning, we have $q'(y|z)$ s.t. $p(y,z)=q'(y|z)p(z)$. Therefore we have:
\begin{align*}
  p(x,y,z)
  &=q(x|z)p(y,z)\\
  &=q(x|z)q'(y|z)p(z)\\
  &=q'(y|z)p(x,z)
\end{align*}

\item (Left) Decomposition: $X, U \Indep Y \; | \; Z \implies U \Indep Y | Z\; $

We are given a $q(x,u|z)$ s.t. $p(x, u, y, z)=q(x, u|z)p(y, z)$.
The result follows by marginalisation over $X$.

\item (Left) Weak Union: $X,U \;\Indep \;  Y \; | \; Z \implies  \; X \Indep Y \; |\; U,Z $

We are given a $q(x,u|z)$ s.t. $p(x, u, y, z)=q(x, u|z)p(y, z)$. By marginalisation over $X$, $p(u, y, z)=q(u|z)p(y, z)$. By conditioning, we have a $q'(x|u,z)$ s.t. $q(x,u|z)=q'(x|u,z)q(u|z)$. Then:
\begin{align*}
  p(x,u,y,z)
  &=q(x,u|z)p(y,z) \\
  &=q'(x|u,z)q(u|z)p(y,z) \\
  &=q'(x|u,z)p(u,y,z)
\end{align*}

\item (Left) Contraction: $(X \Indep Y\; |\; U, Z) \; \land \; (U \Indep Y | Z) \implies X, U \Indep Y \; | \; Z$ \\

We are given $q(x|u,z)$, $q'(u|z)$ s.t.
\begin{align*}
p(x, u, y, z)&=q(x|u,z)p(y, u, z)\\
p(u, y, z)&=q'(u|z)p(y, z)
\end{align*}
We define the product Markov kernel $q''(x,u|z)=q(x|u,z)q'(u|z)$ and have
\begin{align*}
p(x, u, y, z)&=q(x|u,z)q'(u|z)p(y, z)\\
&=q''(x,u|z)p(y, z)
\end{align*}

\item Intersection: For the following part, assume that p is strictly positive, i.e. for all values $x,y,z,u$ we have that $p(x, y, z, u) > 0$. Then show that also the Intersection property holds:
$$(X \Indep Y \;|\;U,Z) \land (X \Indep U \; | \; Y,Z)  \implies  X \Indep U,Y \; | \; Z.$$

We are given $q(x|u,z)$, $q'(x|y,z)$ s.t.
\begin{align*}
p(x, u, y, z)&=q(x|u,z)p(u, y, z)\\
&=q'(x|y,z)p(u,y, z)
\end{align*}
as $p$ is strictly positive, these equalities imply that $\forall x, u, y, z, q(x|u,z)=q'(x|y,z)$. A Markov kernel $q''(x|z)$ thus exists with $\forall x,u,y,z, q''(x|z)=q(x|u,z)=q'(x|y,z)$
\end{enumerate}
\end{gradedsolution}

\section{Reparameterization trick}

\begin{enumerate}
  \item (1pt) Consider the Markov kernel 
    $$K : \RN\times(0,\infty) \dshto \RN : (D,t) \mapsto \int \I_D(x) \C{N}(x \mid t_1,t_2) \,dx$$
    where $\C{N}(x \mid t_1,t_2) = \frac{1}{\sqrt{2\pi t_2^2}} \exp\left(-\frac{(x-t_1)^2}{2 t_2^2}\right)$.
    Derive a deterministic mapping for which this Markov kernel is obtained as the push-forward of the constant uniform Markov kernel $L(U)$ with $\Ucal = [0,1]$.
    
    Hint: You may express your answer in terms of (inverse) error functions (see \href{https://en.wikipedia.org/wiki/Error_function}{wikipedia}).

    Hint: you can use the fact that if two measures on $\RN$ agree on all intervals $[a, b]$, they agree for all elements of the Borel algebra on $\RN$.

    \begin{gradedsolution}
      Define cumulative density function:
      \[
      F_{t_1,t_2} : \RN \to (0, 1) : x \mapsto \int_{-\infty}^x \C N(x \mid t_1, t_2)dx
      \]
      which is strictly monotonic and thus a diffeomorphism with differentiable inverse $F^{-1}_{t_1,t_2}: (0,1) \to \RN$, which is also continuous and thus measurable in $t_1$ and $t_2$ and thus gives a random field:
      \[
        G :(\RN \times (0, \infty)) \times (0,1) \to \RN : (t, u) \mapsto F^{-1}_t(u)
      \]
      giving pushforward Markov kernel (Def 5.2.10):
      \[
        G_*L: \RN \times (0, \infty)  \dshto \RN : (D, t) \mapsto L(U \in (F_t^{-1})^{-1}(D))
      \]
      Let $D$ be a generating set $D=[a,b]$, then:
      \begin{align*}
        U((F_t^{-1})^{-1}(D))
        &= L(U \in \{u \in [0,1] \mid a \le F^{-1}_t(u) \le b \}) \\
        &= L(U \in \{u \in [0,1] \mid F_t(a) \le u \le F_t(b) \}) \\
        &= F_t(b)-F_t(a) \\
        &= \int^b_a \C N(x \mid t_1, t_2)dx\\
      \end{align*}
      which agrees with $K(D|T=t)$. When two measures agree on a generating set of the sigma-algebra, they agree on the entire sigma-algebra.

      From Wikipedia, we see that
      \begin{align*}
        F(x)&=\frac{1}{2}+ \frac{1}{2}\text{erf}(\frac{x-t_1}{\sqrt{2}t_2})\\
        F^{-1}(u)&=\sqrt{2}t_2 \text{erf}^{-1}(2u-1)+t_1
      \end{align*}

    \end{gradedsolution}
  \item (2pts) Same question for the exponential distribution with rate $\lambda > 0$, which has probability density function 
    $$p(x;\lambda) = \I_{[0,\infty)}(x) \lambda e^{-\lambda x}.$$
    First define a corresponding Markov kernel with input $\lambda$. Then derive a deterministic mapping for which that Markov kernel is obtained as the push-forward of the constant uniform Markov kernel $L(U)$ with $\Ucal = [0,1]$.

  \begin{gradedsolution}
    $$K : (0,\infty) \dshto [0,\infty) : (D,\lambda) \mapsto \int_0^\infty \I_D(x) \lambda e^{-\lambda x} \,dx$$
    \begin{align*}
      F_\lambda(b)&=\int_0^b \lambda e^{-\lambda x} dx = 1-e^{-\lambda b}\\
      F_\lambda^{-1}(u)&=\frac{-1}{\lambda}\log(1-u)
    \end{align*}
  \end{gradedsolution}

\end{enumerate}

\end{document}
